\documentclass{article}

\usepackage[utf8]{inputenc}     % pdfLaTeX
\usepackage[T1]{fontenc}
\usepackage[magyar]{babel}

\title{Isaac Newton}
\author{Rovenszky Robin}
\date{Legutóbb frissítve: 2026. Április 16.}

\begin{document}

\maketitle

Isaac Newton a tudományos forradalom időszakában született, amikor a természet jelenségeit már egyre inkább megfigyelésekkel és kísérletekkel próbálták magyarázni. Newton előtt a fizika még nagyrészt a filozófia részének számított, ő azonban döntő szerepet játszott a modern fizika megalapozásában.\\

Newton gyermekkorában nehéz természetű volt, gyakran keveredett konfliktusokba. Tanulmányait a granthami King's Schoolban végezte, majd 1661-ben a Cambridge-i Egyetem Trinity College-ába került. Itt ismerkedett meg azokkal az új tudományos irányzatokkal, amelyek később egész munkásságát meghatározták. A nagy pestisjárvány idején, 1665--1666-ban visszatért szülőfalujába, Woolsthorpe-ba, ahol rendkívül termékeny időszakot élt át; ezt az időszakot gyakran \textit{annus mirabilis}-nek, vagyis csodálatos évnek nevezik.\\

Newton egyik legfontosabb érdeklődési területe a fény volt. Évszázadok óta foglalkoztatta a tudósokat, hogy honnan származnak a színek. Newton egy sötét szobában végzett kísérletek során egy keskeny fénysugarat prizmán vezetett át, és kimutatta, hogy a fehér fény valójában több szín összetétele. Ezzel azt bizonyította, hogy a prizma nem hozza létre a színeket, hanem szétbontja a fényt alkotó összetevőire.\\

E felismerés nyomán Newton arra is rájött, hogy a hagyományos, lencséket használó távcsövek hibáit tükörrel lehet csökkenteni, mivel az üveg a fényt nem egyformán töri meg minden szín esetében. Így megalkotta a tükrös távcső egyik első változatát, a Newton-távcsövet, amely sokkal kisebb színhibával működött, mint a korabeli lencsés eszközök.\\

1669-ben Newton a Cambridge-i Egyetem Lucas-professzora lett. Később visszatért Woolsthorpe-ba, ahol tudományos munkái tovább értek. Ebben az időszakban foglalkozott az analízis alapjaival is; a differenciál- és integrálszámítás kidolgozásában úttörő szerepet játszott, bár a területet később más tudósok, köztük Leibniz is önállóan továbbfejlesztette.\\

Newton legnagyobb hatású művében, a \textit{Philosophi\ae{} Naturalis Principia Mathematica} című könyvben foglalta össze a mozgás törvényeit és az egyetemes gravitáció elméletét. A mű 1687-ben jelent meg, és három alapvető törvényt, valamint a tömegvonzás törvényét tartalmazta. Newton azt állította, hogy minden test vonzza egymást, a vonzóerő pedig a testek tömegétől függ, és a távolság növekedésével gyengül.\\

A gravitációs törvény segítségével Newton rendkívüli pontossággal tudta leírni az égitestek mozgását is. A bolygók pályájáról kimutatta, hogy azok nem kör alakúak, hanem ellipszisek. A gravitáció magyarázatával azt is először tudta tudományosan megindokolni, miért van dagály és apály.\\

Newton három mozgástörvénye a következő:
\begin{enumerate}
    \item \textbf{A tehetetlenség törvénye:} minden test megőrzi nyugalmi állapotát vagy egyenes vonalú egyenletes mozgását, amíg külső erő nem hat rá.
    \item \textbf{A dinamika alaptörvénye:} egy test gyorsulása arányos az eredő erővel, és fordítottan arányos a test tömegével. Mai alakban: $F = m a$.
    \item \textbf{A hatás-ellenhatás törvénye:} minden hatásnak van egy vele egyenlő nagyságú, ellentétes irányú ellenhatása.
\end{enumerate}

Newton nevét Edmund Halley-vel közös munkássága is tovább erősítette. Newton elméletei alapján vizsgálták az üstökösök mozgását, és kimutatták, hogy az egyik üstökös időszakosan, 75 évente visszatér. Tudták, hogy ezt egyikőjük se fogja megélni. Ezt az égitestet ma Halley-üstököseként ismerjük.\\

1693-ban lelki válságon ment keresztül, amelyet később sok életrajzíró Newton ``fekete időszakának'' nevezett. Ennek ellenére munkássága továbbra is meghatározó maradt, és Európa-szerte a korszak egyik leghíresebb tudósaként tartották számon.\\

1703-ban Newton a Royal Society elnöke lett, 1705-ben pedig nemesi címet kapott. Idősebb korában is foglalkoztatta a világ működésének mélyebb megértése, és vallási meggyőződése is erősen befolyásolta gondolkodását. Bár elsősorban tudósként emlékezünk rá, Newton mélyen vallásos volt, és a Biblia jövendölései is fontos szerepet játszottak a világképében.\\

Isaac Newton 1727-ben halt meg, 84 éves korában. Munkássága alapvetően megváltoztatta a természettudományok fejlődését: bebizonyította, hogy a fény, a mozgás és a gravitáció jelenségei egységes törvényekkel írhatók le. A 17. század ezzel a modern tudomány egyik legfontosabb korszakává vált.\\

31 évvel később, karácsony éjszakáján egy csillagász a távcsövével kémlelte az égboltot. És pont ott látott egy égitestet, ahol Halley és Newton előre jósolták az üstököst. Ez erős érv volt Newton törvényei mellett.

\end{document}